\chapter{Arquitectura del sistema}
\label{cap:arquitectura}

\section{Visión general}

La arquitectura de Atlas se ha diseñado para cumplir cuatro objetivos estructurales: desacoplamiento entre servicios, escalado horizontal independiente, tolerancia a fallos por aislamiento de responsabilidades y trazabilidad extremo a extremo de cada operación funcional.

La plataforma se organiza en dos dominios de despliegue:

\begin{itemize}
\item \textbf{Dominio cloud (Azure)}: alberga servicios de exposición pública como API Gateway y CMS.
\item \textbf{Dominio on-premise (CPD municipal)}: ejecuta servicios internos, integración IoT, procesamiento IA y conectividad con sistemas corporativos.
\end{itemize}

Ambos dominios mantienen una arquitectura homogénea basada en contenedores y colas AMQP, permitiendo redistribuir cargas sin afectar al contrato funcional de los microservicios.

\begin{figure}[H]
\centering
\begin{tikzpicture}[
		node distance=1.6cm,
		box/.style={draw, rounded corners, align=center, minimum width=3.6cm, minimum height=1cm},
		flow/.style={->, thick}
	]
	\node[box] (clients) {Clientes web / apps / IoT};
	\node[box, below=of clients] (api) {API Gateway\\(atlas-api)};
	\node[box, below left=of api, xshift=-1.8cm] (broker) {RabbitMQ\\vhost /atlas};
	\node[box, below right=of api, xshift=1.8cm] (cms) {CMS\\(atlas-cms)};
	\node[box, below=of broker] (services) {Microservicios de dominio\\(scraping, IA, conectores)};
	\node[box, right=of services, xshift=1.5cm] (data) {MariaDB / Valkey /\ Ollama};

	\draw[flow] (clients) -- (api);
	\draw[flow] (api) -- node[left] {RPC / F\&F} (broker);
	\draw[flow] (api) -- (cms);
	\draw[flow] (broker) -- (services);
	\draw[flow] (services) -- (data);
	\draw[flow] (cms) -- (data);
\end{tikzpicture}
\caption{Vista lógica de la arquitectura Atlas}
\label{fig:arquitectura-general}
\end{figure}

\section{Patrones de diseño}\label{sec:patrones}
\subsection{Microservicios}\label{sec:microservicios}

La plataforma de ``Smart City'' implementada para el Ayuntamiento de San Lorenzo de el Escorial se basa en un diseño orientado a microservicios, es decir, pequeños procesos dedicados a realizar una tarea específica de manera eficiente.

La nueva Web de Turismo del Ayuntamiento se integra también con la plataforma de ``Smart City'' tal y como se puede ver en \ref{sec:webturismo}

Se disponen de distintos tipos de microservicios:

\begin{itemize}
	\item \textbf{Vistas}: Servicios dedicados a gestionar la interconexión con agentes externos (usuarios, vecinos, dispositivos IoT, \ldots)
	\item \textbf{Controladores}: Servicios orientados a resolver tareas específicas
	\item \textbf{Modelos}: Servicios diseñados para interactuar con sistemas de almacenamiento (Bases de datos, \ldots), APIs externas (tiempo metereológico, estado del tráfico, \ldots), o servicios de Inteligencia Artificial (LLMs, TTS, STT, \ldots)
\end{itemize}

\subsubsection{Microservicios de vista}

Estos microservicios son los encargados de interactuar con el mundo exterior a la plataforma de ``Smart City'':

\paragraph{API Gateway}

Este microservicio expone un API REST para las comunicaciones con vecinos, aplicaciones, \ldots

\paragraph{MQTT}

Sistema de colas compatibles con \textbf{MQTT} (utilizando RabbitMQ) para la recepción de eventos de los dispositivos IoT desplegados a lo largo de la ciudad.

\subsection{Comunicación basada en eventos (RPC y Fire-and-Forget)}

De cara a mantener que todos los microservicios sean capaces de escalar elásticamente en base a las necesidades de cada momento, se diseñan para cumplir dos premisas: \textit{stateless} y \textit{desacoplados}.

Para poder interconectar todos los microservicios se hace uso de eventos y para su gestión descentralizada se dispone de una plataforma de colas basada en el protocolo \textbf{AMQP} (utilizando RabbitMQ).

Dentro de la plataforma \textbf{AMQP} se dispone de dos tipos de eventos:

\begin{itemize}
	\item RPC: Eventos que esperan una respuesta (\textit{Remote Procedure Call}) y por tanto son síncronos.
	\item Fire\&Forget: Eventos que no esperan una respuesta y por tanto son asíncronos.
\end{itemize}

Todos los microservicios implementados en la plataforma hacen uso de la librería común \textit{Ferimer Queue Manager} (\ref{sec:ferimer-queue-manager}), que proporciona una interfaz unificada para productores y consumidores de eventos AMQP.

Esta librería abstrae la complejidad relativa a la gestión de las comunicaciones de la plataforma \textbf{Atlas}.

\subsection{Ontología de mensajería}

La librería `Atlas Messaging Ontology' (\ref{sec:atlas-messaging-ontology}) define todos los posibles mensajes que se transmiten a través de las colas AMQP de la plataforma así como sus esquemas JSON (\textit{JSON Schemas}) utilizados para validar el correcto formato de las peticiones y sus correspondientes respuestas.

\section{Diagrama de flujo de datos}

El flujo principal de datos sigue un patrón de entrada HTTP, normalización de petición, publicación AMQP, ejecución del servicio de dominio y respuesta consolidada. En el caso de eventos asíncronos, la respuesta no bloquea la petición inicial y la trazabilidad se realiza por \textit{correlation id}.

\begin{figure}[H]
\centering
\begin{tikzpicture}[
		node distance=1.8cm,
		proc/.style={draw, rounded corners, align=center, minimum width=3.4cm, minimum height=0.9cm},
		flow/.style={->, thick}
	]
	\node[proc] (client) {Cliente};
	\node[proc, right=of client] (api) {API Gateway};
	\node[proc, right=of api] (broker) {RabbitMQ};
	\node[proc, right=of broker] (service) {Servicio de dominio};
	\node[proc, below=of service] (ia) {Servicio IA / OCR};

	\draw[flow] (client) -- node[above] {HTTP/JSON} (api);
	\draw[flow] (api) -- node[above] {tag + schema} (broker);
	\draw[flow] (broker) -- node[above] {routing key} (service);
	\draw[flow] (service) -- node[right] {invocación} (ia);
	\draw[flow] (service) -- ++(0,-1.5) -| node[below] {respuesta RPC} (api);
	\draw[flow] (api) -- node[below] {HTTP response} (client);
\end{tikzpicture}
\caption{Flujo de datos principal cliente $\rightarrow$ API $\rightarrow$ RabbitMQ $\rightarrow$ servicios}
\label{fig:flujo-datos}
\end{figure}

\section{Seguridad y control de acceso}

La seguridad de la plataforma se implementa por capas, combinando controles de red, validación semántica de mensajes y segmentación de responsabilidades.

\begin{itemize}
\item \textbf{Perímetro y publicación}: NGINX actúa como punto de entrada en ambos entornos y centraliza terminación TLS, redirecciones y cabeceras de seguridad.
\item \textbf{Aislamiento de servicios}: cada microservicio se despliega como pod independiente con permisos mínimos y sin estado local persistente.
\item \textbf{Control de acceso a broker}: RabbitMQ separa cargas por vhosts y credenciales de servicio, limitando permisos de publicación/consumo por cola y exchange.
\item \textbf{Validación de contratos}: todas las peticiones y respuestas AMQP se validan contra \textit{JSON Schemas} de \textit{atlas-messaging-ontology}, reduciendo riesgos de inyección y formatos maliciosos.
\item \textbf{Trazabilidad}: el uso de \textit{correlation id} permite auditoría funcional de cada intercambio RPC y facilita diagnóstico de incidentes.
\end{itemize}

Este enfoque permite alinear la seguridad operativa con el principio de mínimo privilegio y con la naturaleza distribuida de la arquitectura orientada a eventos.
